%%=============================================================================
%% Methodologie
%%=============================================================================

\chapter{Methodologie}%
\label{ch:methodologie}

Dit hoofdstuk beschrijft het plan van aanpak van deze bachelorproef. Het onderzoek
werd opgevat als een iteratief ontwerponderzoek, waarbij inzichten uit literatuur,
doelgroepanalyse en vergelijkende app-analyse vertaald werden naar een prototype,
dat vervolgens kwantitatief en kwalitatief geëvalueerd werd via een survey en
gebruikerstesten.

Het onderzoek werd opgedeeld in vijf opeenvolgende fasen: (1) literatuurstudie,
(2) behoefteanalyse via interviews en app-vergelijking, (3) ontwerp en ontwikkeling
van het prototype, (4) evaluatie via survey en gebruikerstesten, en (5) analyse en
synthese. Per fase worden hieronder de doelstelling, gehanteerde methoden en
gerealiseerde deliverables toegelicht.

\section{Fase 1: Literatuurstudie}

\textbf{Doelstelling.}
In deze fase werd een systematische literatuurstudie uitgevoerd binnen zowel het
probleemdomein als het oplossingsdomein. Enerzijds werd onderzoek bestudeerd rond
ADHD en executieve functies, met bijzondere aandacht voor tijdsblindheid,
initiatieparalyse en uitstelgedrag. Anderzijds werd literatuur geraadpleegd rond
digitale gezondheidstechnologie voor ADHD, cognitieve toegankelijkheid en
UX-ontwerpprincipes voor neurodivergente gebruikers.

\textbf{Methode.}
Wetenschappelijke artikels en relevante vakliteratuur werden verzameld via academische
databanken, waaronder Frontiers in Digital Health, Frontiers in Public Health en
ScienceDirect. Bijzondere aandacht ging naar recente systematische reviews en
meta-analyses van digitale interventies voor ADHD~\autocite{Schofield2026,
Hosseinnia2025}. Naast internationale literatuur werd ook Vlaams onderzoek
geraadpleegd via het onderzoeksproject van HOGENT rond leerstoornissen in het hoger
onderwijs en de UGent-Onderzoeksgroep Ontwikkelingsdiversiteit.

\textbf{Deliverable.}
Een gestructureerde literatuurstudie met synthese van relevante ontwerpprincipes en
een probleemdefiniëring die als theoretische basis diende voor de volgende fasen.

\textbf{Verantwoording.}
Een grondige literatuurstudie is noodzakelijk om het probleem wetenschappelijk te
kaderen en om ontwerpkeuzes te baseren op bestaande inzichten in plaats van
intuïtie.

\section{Fase 2: Behoefteanalyse via interviews en app-vergelijking}

\textbf{Doelstelling.}
Het concretiseren van de centrale onderzoeksvraag vanuit twee complementaire
perspectieven: de ervaring van de doelgroep enerzijds, en het bestaande
app-aanbod anderzijds.

\subsection{Deelfase 2a: Vergelijkende app-analyse}

\textbf{Methode.}
Zeven veelgebruikte planningsapps werden systematisch geanalyseerd op zes
ADHD-relevante criteria: initiatiefdrempel, taakverdeling, dagprioritering, reactie
op gemiste taken, cognitieve belasting en retentie. De geanalyseerde apps waren
Google Calendar, Microsoft To Do, Todoist, Notion, Tiimo, Fabulous en
Goblin.tools. De analyse werd uitgevoerd op basis van directe interactie met de
applicaties en beschikbare documentatie.

\textbf{Deliverable.}
Een vergelijkingstabel met per app een beoordeling op de zes criteria en een
identificatie van de structurele lacune: geen enkele onderzochte app begeleidt
actief het moment ná een gemiste deadline.

\subsection{Deelfase 2b: Semigestructureerde interviews}

\textbf{Methode.}
Er werden twee semigestructureerde, informele interviews afgenomen met studenten
die ADHD-symptomen ervaren. De respondenten werden gerekruteerd via Reddit
(\textit{r/Gent}). De interviews werden afgenomen via Discord en Microsoft Teams
en duurden elk circa 30 tot 45 minuten. Tijdens de gesprekken werd gepeild naar:

\begin{itemize}
  \item huidig planningsgedrag en gehanteerde copingstrategieën;
  \item ervaring met bestaande digitale planningstools;
  \item ervaren moeilijkheden bij het starten en voltooien van taken;
  \item houding tegenover notificaties, autonomie en app-structuur.
\end{itemize}

De interviews werden thematisch geanalyseerd via handmatige codering. Per
respondent werd een gestructureerd verslag opgesteld met thematische samenvatting,
bruikbare citaten en implicaties voor het ontwerp.

\textbf{Deliverable.}
Twee interviewverslagen met thematische analyse en een requirementsdocument
waarin de belangrijkste noden, pijnpunten en functionele vereisten werden
opgelijst.

\textbf{Verantwoording.}
De combinatie van app-analyse en kwalitatieve interviews liet toe om zowel het
bestaande aanbod als de gebruikerservaring systematisch in kaart te brengen.
Beide methoden zijn complementair: de app-analyse identificeert structurele
tekortkomingen op functieniveau, terwijl de interviews de subjectieve beleving
en de psychologische dimensie van planningsproblemen in beeld brengen.

\section{Fase 3: Ontwerp en ontwikkeling van het prototype}

\textbf{Doelstelling.}
Het vertalen van theoretische inzichten, interviewbevindingen en app-analyseresultaten
naar een functioneel, klikbaar prototype dat kan worden ingezet voor
gebruikerstesten.

\subsection{Deelfase 3a: Wireframes}

\textbf{Methode.}
Op basis van de requirements werden drie alternatieve visuele stijlen uitgewerkt,
elk met vier kernschermen: dagweergave, taak toevoegen, gemiste deadline en
patroonherkenning. De drie stijlen verschilden in lay-out en informatiedichtheid:
een kaartgebaseerde stijl (Stijl A), een minimalistische stijl (Stijl B) en een
kleurblokstijl (Stijl C). De wireframes werden gerealiseerd als interactieve
HTML-bestanden.

\textbf{Deliverable.}
Drie vergelijkbare wireframestijlen als downloadbaar HTML- en PDF-bestand.

\subsection{Deelfase 3b: Klikbaar prototype}

\textbf{Methode.}
Op basis van de wireframes werd Stijl A (kaarten) geselecteerd en uitgewerkt tot
een volledig klikbaar prototype met vijf schermen: dagweergave, taakdetail,
taak toevoegen, gemiste deadline en patroonherkenning. Het prototype werd
gerealiseerd als een standalone HTML-bestand zonder backend of
internetverbinding, zodat het eenvoudig kan worden verspreid en uitgevoerd
voor gebruikerstesten. De naam van de applicatie is \textit{Flow}.

De kernfunctionaliteiten van het prototype zijn:
\begin{itemize}
  \item een prikkelarme dagweergave met kleurgecodeerde taakkaarten;
  \item AI-gestuurde taakverdeling in logische stappen met tijdschattingen;
  \item een niet-bestraffende herplanningsflow voor gemiste deadlines;
  \item een patroonherkenningsmodule met automatische planningsaanbevelingen.
\end{itemize}

\textbf{Deliverable.}
Een volledig klikbaar HTML-prototype (\texttt{prototype\_flow\_v2.html}) en een
vereenvoudigde testversie met geïntegreerde taakenlijst voor testpersonen
(\texttt{prototype\_test.html}).

\textbf{Verantwoording.}
De keuze voor een standalone HTML-prototype in plaats van een React- of
React Native-applicatie werd ingegeven door de vereiste van lage distributiedrempel:
testpersonen kunnen het bestand openen via een gewone browser zonder installatie,
account of internetverbinding. Dit verlaagt de drempel voor deelname aan
gebruikerstesten aanzienlijk.

\section{Fase 4: Evaluatie via survey en gebruikerstesten}

\textbf{Doelstelling.}
Het kwantitatief en kwalitatief evalueren van de designkeuzes en de bruikbaarheid
van het prototype binnen de doelgroep.

\subsection{Deelfase 4a: Survey}

\textbf{Methode.}
Een online survey werd opgesteld en verspreid via Google Forms. De survey bestond
uit 23 vragen verdeeld over vijf blokken: profielvragen, vragen over huidige
planningstools, designkeuzes via A/B-comparisons, prototype-evaluatie op basis
van Likert-schalen en een open eindvraag. De survey werd verspreid via de
studentengemeenschap van HOGENT, Reddit (\textit{r/belgium}, \textit{r/ADHD})
en Discord. In totaal namen 32 respondenten deel, waarvan 81,3\% een officiële
diagnose of sterk vermoeden van ADHD had.

De kwantitatieve resultaten werden geanalyseerd op frequentie en procentuele
verdeling. De open antwoorden (n=14) werden thematisch gecodeerd.

\textbf{Deliverable.}
Een surveyvragenlijst, ruwe resultaten en een gestructureerd analyserapport met
bevindingen per blok en ontwerpimplicaties.

\subsection{Deelfase 4b: Gebruikerstesten}

\textbf{Methode.}
Gebruikerstesten werden uitgevoerd met studenten uit de doelgroep. De deelnemers
kregen het prototype (\texttt{prototype\_test.html}) toegestuurd als downloadbaar
bestand en voerden vijf representatieve taken uit: het bekijken van de dagweergave,
het afvinken van een taakstap, het toevoegen van een nieuwe taak, het navigeren
naar de patroonpagina en het kiezen van een herplanoptie op het gemiste
deadline-scherm. De taakenlijst was geïntegreerd in het testbestand zelf, zodat
testpersonen hun voortgang konden bijhouden.

De evaluatie was gericht op:
\begin{itemize}
  \item navigeerbaarheid en begrijpelijkheid van de interface;
  \item ervaren cognitieve belasting;
  \item gepastheid van de herplanningsflow bij gemiste deadlines;
  \item algemene intentie tot gebruik.
\end{itemize}

\textbf{Deliverable.}
Een testverslag met kwalitatieve observaties en aanbevelingen voor verdere
iteratie van het prototype.

\textbf{Verantwoording.}
De combinatie van een kwantitatieve survey en kwalitatieve gebruikerstesten
laat toe om zowel brede tendensen als diepgaande gebruikerservaringen in kaart
te brengen. De survey meet voorkeuren en percepties op schaal; de gebruikerstesten
detecteren concrete interactieproblemen die niet uit zelfrapportage blijken.

\section{Fase 5: Analyse en synthese}

\textbf{Doelstelling.}
Het samenbrengen van alle verzamelde data om conclusies te formuleren per
deelonderzoeksvraag en aanbevelingen op te stellen voor verdere ontwikkeling.

\textbf{Methode.}
De resultaten van de literatuurstudie, app-analyse, interviews, survey en
gebruikerstesten werden geïntegreerd en geïnterpreteerd. Kwalitatieve inzichten
uit de interviews en open survey-antwoorden werden naast de kwantitatieve
surveyresultaten geplaatst om convergerende en divergerende bevindingen te
identificeren. Op basis hiervan werden ontwerpaanbevelingen opgesteld en werd
gereflecteerd over beperkingen en mogelijke vervolgonderzoeken.

\textbf{Deliverable.}
Een geïntegreerde analyse, conclusies per onderzoeksvraag en aanbevelingen voor
verdere ontwikkeling en validatie van de applicatie.

\textbf{Verantwoording.}
Een triangulatie van meerdere databronnen: literatuur, interviews, app-analyse,
survey en gebruikerstesten, verhoogt de validiteit van de conclusies en zorgt
voor een onderbouwde beantwoording van de centrale onderzoeksvraag.

\section{Globale tijdsplanning en deliverables}

De globale planning van het onderzoek verliep als volgt:

\begin{itemize}
  \item \textbf{Week 1--2:} literatuurstudie en vergelijkende app-analyse.
  Deliverable: literatuurmatrix, synthese van ontwerpprincipes en
  vergelijkingstabel van zeven planningsapps op zes ADHD-relevante criteria.

  \item \textbf{Week 2--3:} semigestructureerde interviews en
  requirementsanalyse.
  Deliverable: twee interviewverslagen met thematische analyse en
  requirementsdocument met functionele en niet-functionele vereisten.

  \item \textbf{Week 3--5:} wireframes en eerste prototypeontwikkeling.
  Deliverable: drie wireframestijlen en eerste klikbaar HTML-prototype
  (\textit{Flow} v1).

  \item \textbf{Week 5--6:} survey-opzet, verspreiding en dataverzameling.
  Deliverable: surveyvragenlijst (23 vragen), verspreiding via HOGENT en
  Reddit, 32 ingevulde responsen.

  \item \textbf{Week 6--7:} surveyanalyse en tweede prototype-iteratie.
  Deliverable: analyserapport met bevindingen per blok, en bijgewerkt
  prototype (\textit{Flow} v3) op basis van de surveyresultaten.

  \item \textbf{Week 7--8:} gebruikerstesten, schrijven en aanbevelingen.
  Deliverable: testverslag en eindscriptie met geïntegreerde conclusies
  en ontwerpimplicaties.
\end{itemize}